\documentclass[11pt]{article}
\usepackage[margin=1in]{geometry}
\usepackage{amsmath,amssymb,amsthm}
\usepackage{enumitem}
\usepackage{graphicx}
\usepackage[hidelinks]{hyperref}
\emergencystretch=2em

\newtheorem{theorem}{Theorem}
\newtheorem{corollary}{Corollary}
\newtheorem{remark}{Remark}

\newcommand{\cL}{\mathcal L}
\newcommand{\cK}{\mathcal K}
\newcommand{\NA}{\mathrm{NA}}
\newcommand{\WMP}{\mathrm{WMP}}
\newcommand{\aWMP}{\mathrm{aWMP}}

\title{V-pairs Have the Almost Weak Maximizing Property}
\author{Candidate partial result packet for arXiv:2301.05003}
\date{Run \texttt{fa\_banach\_001}; prepared 2026-07-01}

\begin{document}
\maketitle

\section*{Source Question}

The source paper is Mingu Jung, Gonzalo Mart\'inez-Cervantes, and Abraham
Rueda Zoca, \emph{Rank-one perturbations and norm-attaining operators},
arXiv:2301.05003.  In Question 3.7 the authors ask:
\begin{quote}
Does every \(V\)-pair have the WMP?
\end{quote}
The same page explains why an operator-level approach is insufficient: there
are \(V\)-operators with non-weakly-null maximizing sequences that do not
attain their norms, although the pair in that example is not a \(V\)-pair.

\begin{center}
\includegraphics[width=0.94\linewidth]{figures/open_problem_crop.png}
\end{center}

\section*{Definitions}

A sequence \((x_n)\subset B_X\) is a maximizing sequence for
\(T\in\cL(X,Y)\) if \(\|T x_n\|\to \|T\|\).  A pair \((X,Y)\) has the weak
maximizing property, or \(\WMP\), if every operator \(T:X\to Y\) with a
non-weakly-null maximizing sequence attains its norm.

For this packet, say that \((X,Y)\) has the almost weak maximizing property,
or \(\aWMP\), if every operator \(T:X\to Y\) for which every maximizing
sequence is non-weakly-null attains its norm.  Thus \(\WMP\) implies
\(\aWMP\), but \(\aWMP\) is weaker: it only detects operators whose
maximizing sequences are all non-weakly-null.

We use the standard approximate fixed-point formulation of a \(V\)-operator
used in Khatskevich--Ostrovskii--Shulman and in Remark 3.8 of the source
paper: if \(\|T\|=1\) and \(T\) is a \(V\)-operator, then there are
\(S\in \cL(Y,X)\) with \(\|S\|=1\) and \(y_n\in S_Y\) such that
\[
  \|T S y_n-y_n\|\longrightarrow 0 .
\]
Also, the domain of a \(V\)-pair is reflexive; equivalently, every bounded
sequence in \(X\) has a weakly convergent subsequence.

\section*{Claimed Partial Result}

\begin{theorem}
Every \(V\)-pair has the almost weak maximizing property.
\end{theorem}

\section*{Proof Intuition}

Let \(T:X\to Y\) be a norm-one operator on a \(V\)-pair.  The \(V\)-pair
assumption supplies an almost fixed-point sequence \(y_n\) for \(TS\).  Then
\(S y_n\) is a maximizing sequence for \(T\).  If every maximizing sequence
for \(T\) is non-weakly-null, reflexivity lets us pass to a subsequence
\(S y_n\rightharpoonup x_0\ne0\).  The almost fixed-point relation forces
\(y_n\rightharpoonup T x_0\) and hence \(S T x_0=x_0\).  Since
\(\|S\|=\|T\|=1\), this identity can only happen if
\(\|T x_0\|=\|x_0\|\), so \(T\) attains its norm.

\section*{Proof}

Let \((X,Y)\) be a \(V\)-pair and let \(T\in\cL(X,Y)\) be nonzero.  Normalize
so that \(\|T\|=1\), and assume that every maximizing sequence for \(T\) is
non-weakly-null.

By the \(V\)-pair property and the approximate fixed-point formulation above,
there are \(S\in\cL(Y,X)\), \(\|S\|=1\), and \(y_n\in S_Y\) such that
\[
  \|T S y_n-y_n\|\to 0 .
\]
It follows that
\[
  \|T S y_n\|\to 1 .
\]
Since \(\|S y_n\|\le 1\), the sequence \((S y_n)\subset B_X\) is a maximizing
sequence for \(T\).  By hypothesis, \((S y_n)\) is not weakly null.

The domain \(X\) is reflexive.  Hence, after passing to a subsequence, we may
assume that \(S y_n\rightharpoonup x_0\) for some \(x_0\in B_X\).  Because
the sequence is not weakly null, choose the subsequence so that \(x_0\ne0\).
Then
\[
  T S y_n \rightharpoonup T x_0 .
\]
The norm convergence \(T S y_n-y_n\to0\) gives \(y_n\rightharpoonup T x_0\).
Applying \(S\), we also have \(S y_n\rightharpoonup S T x_0\).  Weak limits
are unique, so
\[
  S T x_0=x_0 .
\]
Therefore
\[
  \|x_0\|=\|S T x_0\|\le \|T x_0\|\le \|x_0\|.
\]
Thus \(\|T x_0\|=\|x_0\|\), and \(x_0\ne0\).  Consequently
\[
  \left\|T\frac{x_0}{\|x_0\|}\right\|=1=\|T\|,
\]
so \(T\) attains its norm.

\section*{Consequence for Question 3.7}

\begin{corollary}
If a \(V\)-pair \((X,Y)\) fails the weak maximizing property, then there is a
non-norm-attaining operator \(T:X\to Y\) with both of the following features:
\begin{enumerate}[label=\textup{(\roman*)}]
\item \(T\) has a non-weakly-null maximizing sequence;
\item \(T\) also has a weakly-null maximizing sequence.
\end{enumerate}
\end{corollary}

\begin{proof}
Failure of \(\WMP\) gives a non-norm-attaining \(T\) with a non-weakly-null
maximizing sequence.  The theorem says that \(T\) cannot have all maximizing
sequences non-weakly-null.  Hence at least one maximizing sequence for \(T\)
is weakly null.
\end{proof}

\section*{Limitations and Novelty Check}

This is only a reduction, not a solution of Question 3.7.  The remaining
problem is exactly the mixed-sequence case in the corollary.

Cheap duplicate checks were run against the run registry, solution, attempt,
and proof-gap indexes for arXiv:2301.05003 and for the core phrases
\(\aWMP\), ``almost weak maximizing'', ``\(V\)-pair'', and ``weak maximizing
property''.  No existing packet for arXiv:2301.05003 or this reduction was
found.  A local later-source search found newer WMP sufficient conditions, in
particular via strict property \((M)\), but did not find an answer to the
exact \(V\)-pair question.  Web searches for the exact phrase ``Does every
V-pair have the WMP'' and nearby variants returned no exact later answer.

Important provenance caution: the arXiv source TeX for arXiv:2301.05003
contains a commented-out draft of essentially the \(\aWMP\) observation.  The
published PDF does not state it.  This packet is therefore best viewed as a
precise promoted cleanup and queue-memory record with modest novelty.

\section*{References}

\begin{enumerate}
\item S. Dantas, M. Jung, and G. Mart\'inez-Cervantes, \emph{Some remarks on
the weak maximizing property}, J. Math. Anal. Appl. 504 (2021), Article
125433.
\item M. Jung, G. Mart\'inez-Cervantes, and A. Rueda Zoca, \emph{Rank-one
perturbations and norm-attaining operators}, arXiv:2301.05003.
\item V. A. Khatskevich, M. I. Ostrovskii, and V. S. Shulman, \emph{Extremal
problems for operators in Banach spaces arising in the study of linear
operator pencils}, Integral Equations Operator Theory 51 (2005), 109--119.
\item M. I. Ostrovskii, \emph{Extremal problems for operators in Banach spaces
arising in the study of linear operator pencils, II}, Integral Equations
Operator Theory 51 (2005), 553--564.
\end{enumerate}

\section*{Human Review Recommendation}

Review as a likely valid partial reduction.  The main convention to verify is
the approximate fixed-point form of \(V\)-operators for the real-space
definition used in the source paper.  Even if accepted, the packet should not
be counted as a full solution to Question 3.7.

\end{document}
